\chapter{Coordinator Agent}

Il Coordinator \`e l'unico agente con visione globale \textit{aggregata}: non pu\`o vedere
tutta la griglia, ma riceve informazioni da tutti gli Scout vicini e costruisce un modello
del mondo pi\`u completo rispetto ai Retriever. Il suo ruolo \`e pianificare chi deve
recuperare cosa, evitando ridondanze e massimizzando l'efficienza del sistema.

\section{Raccolta delle Informazioni}

Ogni volta che riceve un \texttt{ObjectSpottedMessage} da uno Scout, il Coordinator aggiorna
il proprio \texttt{known\_objects} con la posizione e il valore dell'oggetto. Ogni
\texttt{TaskStatusMessage} ricevuto da un Retriever aggiorna uno dei dizionari interni:
\texttt{retriever\_carrying} (numero di oggetti trasportati), \texttt{retriever\_energy}
(energia residua) e \texttt{objects\_being\_collected} (set degli oggetti già in recupero).
Il set \texttt{objects\_being\_collected} \`e la protezione principale contro la doppia
assegnazione: qualsiasi oggetto gi\`a assegnato o dichiarato in recupero da un Retriever
viene escluso dalla lista di candidati per le nuove assegnazioni.

\section{Logica di Assegnazione}

La funzione \texttt{\_plan\_task\_assignments} calcola per ogni coppia (Retriever, oggetto)
una priorit\`a greedy:

$$\text{priority}(o, r) = \frac{\text{value}(o)}{d(r, o) + 1} + b_{\text{spotted}}$$

dove $b_{\text{spotted}} = 0{,}5$ se il Retriever ha gi\`a visto l'oggetto direttamente
(campo \texttt{\_spotted\_by}), altrimenti $0$. I candidati vengono ordinati per priorit\`a
decrescente e assegnati in modo greedy rispettando i slot liberi di ciascun Retriever.
Al momento dell'invio, il Coordinator trasmette sia il \texttt{TaskAssignmentMessage}
che un \texttt{MapDataMessage} con le celle magazzino note, accelerando la scoperta delle
stazioni di deposito da parte dei Retriever che non le hanno ancora esplorate.

\section{Riposizionamento verso il Centroide}

Quando il Coordinator non ha compiti da assegnare, il suo obiettivo di movimento \`e
rimanere raggiungibile dalla propria flotta. All'avvio, finch\'e non riceve nessuna
posizione recente da un Retriever, il Coordinator percorre un circuito di cinque waypoint
strategici (centro e quattro quadranti della griglia, agganciati alla cella libera pi\`u
vicina) per coprire tutta la griglia e ristabilire il contatto nel minor numero di passi.

Non appena arrivano posizioni fresche (pi\`u recenti di 25 passi), calcola il centroide
medio di tutti i Retriever e, se la propria distanza da quel punto supera una soglia
(pari a $\frac{1}{4}$ del raggio di comunicazione), si sposta verso la cella libera pi\`u
vicina al centroide.

Se la distanza \`e gi\`a entro la soglia, il Coordinator rimane fermo (\texttt{IDLE}):
non esegue n\`e esplorazione frontier-based n\`e passeggiate casuali, poich\`e il suo
compito non \`e esplorare ma coordinarsi. Questa scelta riduce il consumo energetico e
evita che l'agente si allontani involontariamente dagli agenti che gestisce.

Se esistono oggetti noti ma nessun Retriever \`e nel raggio di comunicazione, il Coordinator
entra in modalit\`a \textbf{SEEK-RETRIEVER}: si dirige verso l'ultima posizione dichiarata
di un Retriever noto, con l'obiettivo di ristabilire il canale di comunicazione e procedere
con l'assegnazione appena possibile.

\section{Sincronizzazione tra Coordinator Multipli}

In configurazioni con pi\`u di un Coordinator, il rischio di doppia assegnazione aumenta.
Il \texttt{CoordinatorSyncMessage} viene scambiato periodicamente tra Coordinator vicini
e contiene la lista delle assegnazioni attive (\texttt{assigned\_tasks}) e degli oggetti
in recupero (\texttt{objects\_being\_collected}). Alla ricezione, ogni Coordinator fonde
le informazioni ricevute con le proprie, aggiornando i set interni. Questa sincronizzazione
non \`e istantanea — si propaga a onde tramite la comunicazione locale — ma \`e sufficiente
per eliminare la stragrande maggioranza delle doppie assegnazioni in scenari realistici.
